\section{Introduction}

Feynman's parton distribution functions (PDFs) \cite{Feynman:1973xc}  $f(x)$
are  the crucial building blocks in the description of hard inclusive processes
in quantum chromodynamics (QCD).
Accumulating nonperturbative  information about
the hadron structure,   the PDFs  are a natural subject for
a lattice study.
 However, straightforward definitions of  PDFs
refer to    matrix elements
of bilocal operators on   the light cone $z^2=0$,
the intervals inaccessible on   the   Euclidean lattice.

The ideas of how to    get  information  from space-like intervals
date to the pioneering paper of  \mbox{W. Detmold}   and D. Lin
\cite{Detmold:2005gg}  who proposed a   lattice study
of  the deep-inelastic-type Euclidean correlators   of heavy-light  currents.
Later, V. Braun and D. M\"uller
\cite{Braun:2007wv}
proposed to use   Euclidean correlators to
extract  the pion
distribution amplitude,
another   function \cite{Radyushkin:1977gp} playing a fundamental role in  perturbative QCD studies
of  hard exclusive processes.
The use of correlators in the form of ``lattice cross sections''  was more recently
  advocated in the papers
by Qiu and collaborators \cite{Ma:2014jla,Ma:2017pxb}.

The current correlators involve a quark propagator connecting the current vertices.
This factor is avoided in   the  proposal   by X. Ji   \cite{Ji:2013dva}  to study  the
quasi-PDFs $Q(y,p_3)$  that describe the  distribution
of the spatial $z_3$-component of the hadron momentum $p_3$.
While being different from the Feynman PDFs $f(y)$ describing the distribution
of the hadron's ``plus''-momentum $p_+=p_0+p_3$,
they  coincide with  $f(y)$ in  the infinite momentum limit
$p_3\to \infty$.

Both on the lattice and in the usual continuum space,
the basic object for all types of PDFs is the
matrix element $M(z,p)$ generically
(i.e. ignoring the inessential spin complications)
 written as $\langle  p | \phi (0)  \phi(z)  | p \rangle $.
 By Lorentz invariance, it is a function of the
 Ioffe time $(pz)\equiv -\nu$ \cite{Ioffe:1969kf} and the interval $z^2$,
 $M(z,p) \equiv {\cal M} (\nu, -z^2)$.

 In the (formal)  light-cone  limit $z^2=0$, the Fourier transform
 of  $ {\cal M} (\nu, 0)$ with respect to $\nu$
 gives $f(x)$.  In this sense, the $\nu$-dependence of
 the Ioffe-time distribution (ITD) ${\cal M} (\nu, -z^2)$ reflects
 the longitudinal structure of the PDFs.
 As shown in \mbox{Ref. \cite{Radyushkin:2016hsy}},
 the \mbox{$z^2$-dependence}  of  ${\cal M} (\nu, -z^2)$ determines
 the \mbox{$k_\perp$-dependence}  of the  (straight-link in the case of QCD)
  transverse momentum dependent
 parton distributions (TMDs) ${\cal F} (x, k_\perp)$.


 Since the quasi-PDFs $Q(y,p_3)$ are given by
 the Fourier transform of ${\cal M} (z_3 p_3, z_3^2)$
 with respect to $z_3$, their shape  is distorted by nonperturbative
 transverse momentum effects  entering through the second argument of
  ${\cal M} (z_3 p_3, z_3^2)$.
 While,   in a general perspective, the
 \mbox{$k_\perp$-dependence}  of  ${\cal F} (x, k_\perp)$
provides information about the three-dimensional structure
of hadrons, in the case of the quasi-PDFs it
is a nuisance responsible for the unwanted  difference between
$Q(y,p_3)$ and $f(y)$ that is very strong at momenta reached in
existing lattice calculations of quasi-PDFs.

To decrease the impact of the $z^2$-dependence of the
ITD ${\cal M} (\nu, -z^2)$, it was proposed  \cite{Radyushkin:2017cyf}
to consider the reduced ITD ${\mathfrak M} (\nu, -z^2)$
given by the ratio of  ${\cal M} (\nu, -z^2)$ and the rest-frame
distribution  ${\cal M} (0, -z^2)$.
Though  there are no first-principle grounds
that the nonperturbative part of the \mbox{$z^2$-dependence}  disappears in this ratio,
it is natural to expect that it is strongly reduced.

The ideal case when ${\mathfrak M} (\nu, -z^2)$ is  just a function
of $\nu$  corresponds to factorization
of the $x$ and $k_\perp$ dependencies of the TMD
${\cal F} (x, k_\perp)$.  In fact, the idea that ${\cal F} (x, k_\perp) = f(x) K(k_\perp)$
in the soft region $k_\perp^2 \lesssim 1$ GeV$^2$  is
a standard assumption of the TMD practitioners  (see, e.g., Ref. \cite{Anselmino:2013lza}),
with a Gaussian being the most popular form for
$K(k_\perp)$.

An exploratory lattice study of the reduced ITD
was  performed in Ref. \cite{Orginos:2017kos} (and also described in Ref.
\cite{Karpie:2017bzm}).  The results show that
${\mathfrak M} (\nu, z_3^2)$ is basically a universal
function of $\nu$, with small deviations from the common curve
for the points corresponding to the smallest values of $z_3$.

As demonstrated in Ref. \cite{Orginos:2017kos}, these
deviations may be explained by  perturbative evolution.
While  the leading logarithmic approximation (LLA)
used in Ref.  \cite{Orginos:2017kos}  is  sufficient to analyze the
$\ln z_3^2$ dependence,
one needs to go   beyond it to
specify the scale $\mu$  which should be attributed
to the extracted scale-dependent PDFs $f(x,\mu^2)$.



To this end,  one needs complete expressions for one-loop corrections to ITDs.
Recently,  such  calculations have been reported  in Refs.
 \cite{Ji:2017rah,Radyushkin:2017lvu}.
Our goal here is to give a more detailed discussion
of the LLA  treatment of the evolution,
and also to extend the  analysis
beyond the LLA.
As we will  show,  the one-loop correction contains a large
 contribution  that  considerably changes the results
 obtained in the LLA.


 To make this article  self-contained, we outline  in \mbox{Sec. II}
 the basics of  the Ioffe-time distributions and pseudo-PDFs.
 In Sec.  III, we discuss the structure of  one-loop corrections.
 In Sec.  IV, we  describe the evolution effects
 revealed  in the lattice study of Ref. \cite{Orginos:2017kos},
 and convert the data for  the reduced ITD ${\mathfrak M} (\nu, z_3^2)$
 into the standard parton densities $f(x,\mu^2)$ defined in the $\overline{\rm MS}$  scheme.
 The summary of the paper and conclusions  are  given in \mbox{Sec.  V.}
